%==============================================================================
% 1. INTRODUCTION
%==============================================================================
% Drafted from paper/.research/{motivation,landscape,results}.md.
% Cite keys are the ones defined in refs.bib:
%   maharana2024locomo, wu2024longmemeval, beam2025, mem0_2024, zep_graphiti_2025
% (the brief names locomo/longmemeval/beam/mem0/zep map onto these).
%------------------------------------------------------------------------------
\section{Introduction}
\label{sec:intro}

When a company decides to raise enterprise pricing, the useful artifact is not the
new number. It is the record that the number was raised from a prior value, because
of churn observed in a specific customer cohort, reviewed on a specific date by a
specific person, with a pending re-evaluation triggered by a named competitor's
move. An AI agent asked ``when did we change enterprise pricing, and why?'' must
recover not the current state but the transition. This is \emph{organizational
memory}: the accumulated record of what a company
decided, why, when, and who was involved, scattered across every surface where work
happens, including Slack threads, meeting transcripts, tickets, emails, shared
documents, and contracts. It is the substrate that lets independent agents converge
on one consistent account of what is true, what was true, and why it changed.
Without it, each agent carries its own partial context and acts on facts the rest of
the organization has already superseded; the company gets \emph{less} coherent the
more agents are added to it. As agents move from individual productivity tools to
operational participants doing contract review, customer escalation, engineering
planning, and financial modeling, the binding failure mode is no longer ``the agent
does not know.'' It is ``the agent confidently acts on a fact that is no longer
true.''

Organizational memory is structurally unlike the personal, conversational memory
that current benchmarks measure, in three ways that jointly make it hard. It is
\emph{graph-shaped}: the relationships between facts carry as much weight as the
facts themselves, because a decision is only actionable alongside the conversation
that prompted it, the estimate that bounded it, and the sign-off that ratified it.
It is \emph{bi-temporal}: companies revise decisions, and a usable memory must date
and recover both the change and the prior state independently, separating what was
true \emph{then} (valid time) from when the system learned it (ingestion time), so
that supersession is preserved rather than overwritten. And it is
\emph{decision-traced}: the reasoning behind a fact is frequently more valuable than
the fact, because decision traces are how organizations avoid re-litigating settled
questions. A corpus with these properties stops looking like a chat log and starts
looking like a company: dozens of simultaneous authors writing across incompatible
channels, recording contradictions across teams that the system must surface rather
than silently resolve.

The dominant memory benchmarks do not exercise any of this. LoCoMo
\citep{maharana2024locomo} evaluates very-long-term \emph{two-person} dialogue and is
now near-saturated; LongMemEval \citep{wu2024longmemeval} tests knowledge updates and
temporal reasoning but only over a \emph{single user's own stated facts}; and BEAM
\citep{beam2025} pushes the volume frontier to coherent \emph{single-narrator}
conversation that, though unsaturated, still asks \emph{how much} one person said, not
\emph{who} said what across an organization (per-benchmark scale and SOTA figures, and
the disputed-LoCoMo argument, in \S\ref{sec:related:benchmarks}). All three share one
structural assumption: memory is a single narrator's history, possibly very long. None
varies author count, channel heterogeneity, role and scope structure, or attributed
supersession, which are precisely the axes that define organizational memory. The
near-saturation of LoCoMo and LongMemEval is therefore not evidence that memory is
solved; it is evidence that these benchmarks stopped measuring the hard part.

\bench{} targets exactly the regime these benchmarks omit. It is a benchmark of
long-horizon organizational memory built over a multi-channel synthetic corpus that
behaves like a real company's record: an eighty-five-person organization writing
across email, Slack, meeting transcripts, tickets, wikis, and CRM, over a multi-year
history in which facts are routinely superseded, corrected, and left in unresolved
contradiction. Its questions are not retrievals from a narrator's history but
cross-channel, multi-hop inferences over years of accumulated artifact, organized
into a taxonomy of organizational competencies. The small and medium tiers evaluated
here exercise six of these competencies (C1--C6): supersession, decision provenance,
bi-temporal (as-of-date) reasoning, audit replay, justification chains, and
contradiction resolution. A seventh competency, emergent patterns assembled from
behavior that was never written down as a rule, is evaluated only at the large tier
and above, where a deep enough behavioural corpus exists for the patterns to be
reconstructable; it is out of scope for the small and medium tiers in this paper.
We evaluate four publicly available knowledge-graph and memory systems that we
could obtain and run end-to-end against an external corpus: gbrain, Zep/Graphiti
\citep{zep_graphiti_2025}, the hosted Mem0 platform \citep{mem0_2024}, and the
open-source graphify. All four run under a single fixed harness with an identical
answerer and judge, and the headline result is not that any system wins:
Table~\ref{tab:leaderboard} shows that no system is close, the strongest public
system, gbrain, reaching only $0.394$ small and $0.428$ medium. The small-tier
ordering does not even survive the move to medium scale, so conversational-scale
behavior does not predict organizational-scale behavior, and by this measurement
organizational-scale memory is an \emph{unsolved} problem with large remaining
headroom.

\paragraph{Contributions.}
\begin{itemize}
  \item \textbf{The \bench{} benchmark.} A benchmark for long-horizon organizational
        memory over a multi-channel, multi-author, multi-year synthetic corpus
        (small and medium tiers of $121$ and $443$ artifacts, scaling toward larger
        tiers), with questions spanning organizational competencies that no
        prior conversational benchmark exposes. The small and medium tiers evaluated
        here cover six competencies (C1--C6; Table~\ref{tab:categories}); a seventh,
        emergent patterns, is evaluated only at the large tier and above and is out
        of scope here.
  \item \textbf{A reproducible, neutral evaluation harness} that runs the four
        publicly available memory systems (gbrain, Zep/Graphiti, Mem0 platform, and
        open-source graphify) under identical conditions: shared answerer, shared
        LLM-as-judge with a fixed rubric, and a \emph{readiness-gating} methodology
        that drains asynchronous hosted systems on a probe search rather than a naive
        ingest signal, a fix that materially changes measured scores and that we
        document as a contribution in its own right.
  \item \textbf{The empirical finding that the field is unsolved.} No publicly
        available system comes close to organizational competence: every system
        collapses on justification chains, bi-temporal reasoning, and contradiction,
        and the ranking among the trailing systems reorders between tiers, indicating
        a ceiling far below usable performance and large headroom for future work
        (full per-tier, per-category results in Section~\ref{sec:experiments}).
  \item \textbf{Public release.} We release the corpus, the question set, the
        evaluation harness, and a leaderboard, so that subsequent systems can be
        scored under the same fixed conditions.
\end{itemize}

The remainder of the paper situates \bench{} against prior memory benchmarks and
systems (Section~\ref{sec:related}), details the corpus and question-generation
design (Section~\ref{sec:design}), describes the baseline systems and their
configuration (Section~\ref{sec:baselines}), reports and analyzes results
(Section~\ref{sec:experiments}), and closes with limitations
(Section~\ref{sec:limitations}) and broader implications.

% TEASER FIGURE --- Figure 1, in the intro (guide 2.7 / 5.2). Guarded with
% \IfFileExists so the section compiles before the asset exists.
\begin{figure}[t]
  \centering
  \resizebox{0.95\linewidth}{!}{% =============================================================================
% figures/teaser_diagram.tex --- OrgMemBench Figure 1 teaser / system diagram.
%
% RATIONALE (what this changes vs. the old placeholder, and why).
% The placeholder was a flat 4-box chain (corpus -> questions -> system ->
% eval) that hid the one idea that makes OrgMemBench distinctive: a single
% bi-temporal SOURCE-OF-TRUTH GRAPH is the origin of *both* the multi-channel
% corpus *and* the questions with their gold answers. This figure makes that
% DUAL PROJECTION the visual anchor: one graph node fans out to (a) the
% multi-channel synthetic corpus and (b) the competency-tagged questions with
% deterministic graph-projected gold answers. It then shows the real eval loop
% the chain elided: each corpus is ingested by the four systems under test
% (gbrain / zep-cloud / mem0-platform / graphify-oss), retrieval feeds a shared
% (or native) answerer, and a rubric-based, capability-aware judge scores the
% answer *against the graph-projected gold* -> per-category rubric mean ->
% leaderboard. The six competencies C1-C6 ride along as tags on the questions,
% and the gold answers visibly trace back to the same graph the judge scores
% against (the dashed "scored against" link), which is the property the flat
% chain could not express. Restrained: single orange accent, the rest grey.
%
% PREAMBLE main.tex needs (add once):
%   \usepackage{tikz}
%   \usetikzlibrary{positioning,arrows.meta,shapes.geometric,fit,backgrounds}
% DROP-IN for the fig:teaser block (replace the \IfFileExists/\includegraphics):
%   % =============================================================================
% figures/teaser_diagram.tex --- OrgMemBench Figure 1 teaser / system diagram.
%
% RATIONALE (what this changes vs. the old placeholder, and why).
% The placeholder was a flat 4-box chain (corpus -> questions -> system ->
% eval) that hid the one idea that makes OrgMemBench distinctive: a single
% bi-temporal SOURCE-OF-TRUTH GRAPH is the origin of *both* the multi-channel
% corpus *and* the questions with their gold answers. This figure makes that
% DUAL PROJECTION the visual anchor: one graph node fans out to (a) the
% multi-channel synthetic corpus and (b) the competency-tagged questions with
% deterministic graph-projected gold answers. It then shows the real eval loop
% the chain elided: each corpus is ingested by the four systems under test
% (gbrain / zep-cloud / mem0-platform / graphify-oss), retrieval feeds a shared
% (or native) answerer, and a rubric-based, capability-aware judge scores the
% answer *against the graph-projected gold* -> per-category rubric mean ->
% leaderboard. The six competencies C1-C6 ride along as tags on the questions,
% and the gold answers visibly trace back to the same graph the judge scores
% against (the dashed "scored against" link), which is the property the flat
% chain could not express. Restrained: single orange accent, the rest grey.
%
% PREAMBLE main.tex needs (add once):
%   \usepackage{tikz}
%   \usetikzlibrary{positioning,arrows.meta,shapes.geometric,fit,backgrounds}
% DROP-IN for the fig:teaser block (replace the \IfFileExists/\includegraphics):
%   % =============================================================================
% figures/teaser_diagram.tex --- OrgMemBench Figure 1 teaser / system diagram.
%
% RATIONALE (what this changes vs. the old placeholder, and why).
% The placeholder was a flat 4-box chain (corpus -> questions -> system ->
% eval) that hid the one idea that makes OrgMemBench distinctive: a single
% bi-temporal SOURCE-OF-TRUTH GRAPH is the origin of *both* the multi-channel
% corpus *and* the questions with their gold answers. This figure makes that
% DUAL PROJECTION the visual anchor: one graph node fans out to (a) the
% multi-channel synthetic corpus and (b) the competency-tagged questions with
% deterministic graph-projected gold answers. It then shows the real eval loop
% the chain elided: each corpus is ingested by the four systems under test
% (gbrain / zep-cloud / mem0-platform / graphify-oss), retrieval feeds a shared
% (or native) answerer, and a rubric-based, capability-aware judge scores the
% answer *against the graph-projected gold* -> per-category rubric mean ->
% leaderboard. The six competencies C1-C6 ride along as tags on the questions,
% and the gold answers visibly trace back to the same graph the judge scores
% against (the dashed "scored against" link), which is the property the flat
% chain could not express. Restrained: single orange accent, the rest grey.
%
% PREAMBLE main.tex needs (add once):
%   \usepackage{tikz}
%   \usetikzlibrary{positioning,arrows.meta,shapes.geometric,fit,backgrounds}
% DROP-IN for the fig:teaser block (replace the \IfFileExists/\includegraphics):
%   \input{figures/teaser_diagram}
% Keep the existing \caption{...}\label{fig:teaser}; do not caption inside here.
% =============================================================================
\begin{tikzpicture}[
    font=\small,
    node distance=6mm,
    accent/.style={orange!75!red},
    % --- the bi-temporal source-of-truth graph (origin) ---
    graph/.style={cylinder, shape border rotate=90, aspect=0.32, draw=orange!75!red,
                  very thick, fill=orange!8, minimum width=20mm, minimum height=15mm,
                  align=center, inner sep=2pt, text=black},
    % --- generic stage box ---
    box/.style={draw=black!55, thick, rounded corners=2pt, fill=black!2,
                align=center, inner sep=4pt, minimum height=9mm, text width=24mm},
    % --- system-under-test chip ---
    sys/.style={draw=black!45, semithick, rounded corners=1.5pt, fill=white,
                font=\scriptsize\ttfamily, inner sep=2.5pt, minimum width=23mm,
                align=center},
    % --- competency tag ---
    tag/.style={draw=orange!70!red, semithick, rounded corners=1pt, fill=orange!10,
                font=\scriptsize, inner sep=1.6pt, text=black!85},
    % --- judge / output emphasis ---
    judge/.style={draw=orange!75!red, very thick, rounded corners=2pt, fill=orange!6,
                  align=center, inner sep=4pt, minimum height=9mm, text width=26mm},
    board/.style={draw=black!60, thick, rounded corners=2pt, fill=black!4,
                  align=center, inner sep=4pt, minimum height=9mm, text width=22mm},
    flow/.style={-{Latex[length=2mm,width=1.6mm]}, black!65, semithick},
    proj/.style={-{Latex[length=2mm,width=1.6mm]}, orange!75!red, thick},
    scored/.style={{Latex[length=1.6mm]}-{Latex[length=1.6mm]}, densely dashed,
                   orange!75!red, semithick},
  ]

  % ---- ORIGIN: bi-temporal source-of-truth graph -------------------------
  \node[graph] (graph) {\textbf{Source-of-truth}\\\textbf{graph}\\[1pt]
        \scriptsize\itshape bi-temporal\\[-1pt]
        \scriptsize\ttfamily valid/invalid\\[-2pt]
        \scriptsize\ttfamily ingested/expired};
  \node[font=\scriptsize\itshape, text=black!55, above=1mm of graph]
        {synthetic Helix org};

  % ---- DUAL PROJECTION targets (the key idea) ----------------------------
  \node[box, above right=6mm and 16mm of graph, text width=27mm] (corpus)
        {\textbf{Multi-channel corpus}\\[1pt]
         \scriptsize email $\cdot$ Slack $\cdot$ transcripts\\[-2pt]
         \scriptsize tickets $\cdot$ docs $\cdot$ CRM};
  \node[box, below right=6mm and 16mm of graph, text width=27mm] (gold)
        {\textbf{Questions + gold}\\[1pt]
         \scriptsize competency-tagged Qs\\[-2pt]
         \scriptsize deterministic gold answers};

  % competency tags C1-C6 hanging under the questions
  \node[tag, below=2.5mm of gold] (tags)
        {C1 supers.\ $\cdot$ C2 prov.\ $\cdot$ C3 bi-temp.\ $\cdot$
         C4 audit $\cdot$ C5 justif.\ $\cdot$ C6 contra.};

  % projection arrows from the one graph (orange = "projected from graph")
  \draw[proj] (graph.east) to[out=35,in=180]
        node[pos=0.55, above=-0.5mm, accent, font=\scriptsize\itshape] {project}
        (corpus.west);
  \draw[proj] (graph.east) to[out=-35,in=180] (gold.west);

  % ---- SYSTEMS UNDER TEST -------------------------------------------------
  \node[box, right=14mm of corpus, text width=15mm, fill=white, draw=black!45]
        (ingest) {\textbf{ingest}\\\scriptsize+ retrieve};
  \node[sys, above right=2mm and 9mm of ingest] (s1) {gbrain};
  \node[sys, below=1.2mm of s1] (s2) {zep-cloud};
  \node[sys, below=1.2mm of s2] (s3) {mem0-platform};
  \node[sys, below=1.2mm of s3] (s4) {graphify-oss};
  \begin{scope}[on background layer]
    \node[draw=black!30, rounded corners=2pt, fill=black!2, fit=(s1)(s4),
          inner sep=2.5pt] (systems) {};
  \end{scope}
  \node[font=\scriptsize\itshape, text=black!55, above=0.5mm of systems]
        {memory systems under test};

  \draw[flow] (corpus.east) -- (ingest.west);
  \draw[flow] (ingest.east) -- (systems.west);

  % ---- ANSWERER -----------------------------------------------------------
  \node[box, right=10mm of systems, text width=20mm] (answerer)
        {\textbf{Answerer}\\[1pt]
         \scriptsize shared neutral\\[-2pt]
         \scriptsize (or native synth.)};
  \draw[flow] (systems.east) -- (answerer.west);

  % ---- JUDGE (capability-aware, rubric) ----------------------------------
  \node[judge, below=12mm of answerer] (judge)
        {\textbf{Capability-aware judge}\\[1pt]
         \scriptsize rubric facet coverage\\[-2pt]
         \scriptsize readiness-gated};
  \draw[flow] (answerer.south) -- node[right=0.5mm, font=\scriptsize, text=black!60]
        {answer} (judge.north);

  % ---- LEADERBOARD --------------------------------------------------------
  \node[board, left=12mm of judge] (board)
        {\textbf{Leaderboard}\\[1pt]
         \scriptsize per-category\\[-2pt]
         \scriptsize rubric mean};
  \draw[flow] (judge.west) -- (board.east);

  % gold answers feed the judge, and visibly trace back to the SAME graph
  \draw[flow] (gold.east) to[out=0,in=205]
        node[pos=0.5, below=0.3mm, font=\scriptsize, text=black!60] {gold}
        (judge.west);
  % the dashed "scored against" link: gold (graph-projected) <-> judged answer
  \draw[scored] (graph.south) to[out=-90,in=180]
        node[pos=0.7, below=0.5mm, accent, font=\scriptsize\itshape]
        {scored against graph truth} (board.west);

\end{tikzpicture}

% Keep the existing \caption{...}\label{fig:teaser}; do not caption inside here.
% =============================================================================
\begin{tikzpicture}[
    font=\small,
    node distance=6mm,
    accent/.style={orange!75!red},
    % --- the bi-temporal source-of-truth graph (origin) ---
    graph/.style={cylinder, shape border rotate=90, aspect=0.32, draw=orange!75!red,
                  very thick, fill=orange!8, minimum width=20mm, minimum height=15mm,
                  align=center, inner sep=2pt, text=black},
    % --- generic stage box ---
    box/.style={draw=black!55, thick, rounded corners=2pt, fill=black!2,
                align=center, inner sep=4pt, minimum height=9mm, text width=24mm},
    % --- system-under-test chip ---
    sys/.style={draw=black!45, semithick, rounded corners=1.5pt, fill=white,
                font=\scriptsize\ttfamily, inner sep=2.5pt, minimum width=23mm,
                align=center},
    % --- competency tag ---
    tag/.style={draw=orange!70!red, semithick, rounded corners=1pt, fill=orange!10,
                font=\scriptsize, inner sep=1.6pt, text=black!85},
    % --- judge / output emphasis ---
    judge/.style={draw=orange!75!red, very thick, rounded corners=2pt, fill=orange!6,
                  align=center, inner sep=4pt, minimum height=9mm, text width=26mm},
    board/.style={draw=black!60, thick, rounded corners=2pt, fill=black!4,
                  align=center, inner sep=4pt, minimum height=9mm, text width=22mm},
    flow/.style={-{Latex[length=2mm,width=1.6mm]}, black!65, semithick},
    proj/.style={-{Latex[length=2mm,width=1.6mm]}, orange!75!red, thick},
    scored/.style={{Latex[length=1.6mm]}-{Latex[length=1.6mm]}, densely dashed,
                   orange!75!red, semithick},
  ]

  % ---- ORIGIN: bi-temporal source-of-truth graph -------------------------
  \node[graph] (graph) {\textbf{Source-of-truth}\\\textbf{graph}\\[1pt]
        \scriptsize\itshape bi-temporal\\[-1pt]
        \scriptsize\ttfamily valid/invalid\\[-2pt]
        \scriptsize\ttfamily ingested/expired};
  \node[font=\scriptsize\itshape, text=black!55, above=1mm of graph]
        {synthetic Helix org};

  % ---- DUAL PROJECTION targets (the key idea) ----------------------------
  \node[box, above right=6mm and 16mm of graph, text width=27mm] (corpus)
        {\textbf{Multi-channel corpus}\\[1pt]
         \scriptsize email $\cdot$ Slack $\cdot$ transcripts\\[-2pt]
         \scriptsize tickets $\cdot$ docs $\cdot$ CRM};
  \node[box, below right=6mm and 16mm of graph, text width=27mm] (gold)
        {\textbf{Questions + gold}\\[1pt]
         \scriptsize competency-tagged Qs\\[-2pt]
         \scriptsize deterministic gold answers};

  % competency tags C1-C6 hanging under the questions
  \node[tag, below=2.5mm of gold] (tags)
        {C1 supers.\ $\cdot$ C2 prov.\ $\cdot$ C3 bi-temp.\ $\cdot$
         C4 audit $\cdot$ C5 justif.\ $\cdot$ C6 contra.};

  % projection arrows from the one graph (orange = "projected from graph")
  \draw[proj] (graph.east) to[out=35,in=180]
        node[pos=0.55, above=-0.5mm, accent, font=\scriptsize\itshape] {project}
        (corpus.west);
  \draw[proj] (graph.east) to[out=-35,in=180] (gold.west);

  % ---- SYSTEMS UNDER TEST -------------------------------------------------
  \node[box, right=14mm of corpus, text width=15mm, fill=white, draw=black!45]
        (ingest) {\textbf{ingest}\\\scriptsize+ retrieve};
  \node[sys, above right=2mm and 9mm of ingest] (s1) {gbrain};
  \node[sys, below=1.2mm of s1] (s2) {zep-cloud};
  \node[sys, below=1.2mm of s2] (s3) {mem0-platform};
  \node[sys, below=1.2mm of s3] (s4) {graphify-oss};
  \begin{scope}[on background layer]
    \node[draw=black!30, rounded corners=2pt, fill=black!2, fit=(s1)(s4),
          inner sep=2.5pt] (systems) {};
  \end{scope}
  \node[font=\scriptsize\itshape, text=black!55, above=0.5mm of systems]
        {memory systems under test};

  \draw[flow] (corpus.east) -- (ingest.west);
  \draw[flow] (ingest.east) -- (systems.west);

  % ---- ANSWERER -----------------------------------------------------------
  \node[box, right=10mm of systems, text width=20mm] (answerer)
        {\textbf{Answerer}\\[1pt]
         \scriptsize shared neutral\\[-2pt]
         \scriptsize (or native synth.)};
  \draw[flow] (systems.east) -- (answerer.west);

  % ---- JUDGE (capability-aware, rubric) ----------------------------------
  \node[judge, below=12mm of answerer] (judge)
        {\textbf{Capability-aware judge}\\[1pt]
         \scriptsize rubric facet coverage\\[-2pt]
         \scriptsize readiness-gated};
  \draw[flow] (answerer.south) -- node[right=0.5mm, font=\scriptsize, text=black!60]
        {answer} (judge.north);

  % ---- LEADERBOARD --------------------------------------------------------
  \node[board, left=12mm of judge] (board)
        {\textbf{Leaderboard}\\[1pt]
         \scriptsize per-category\\[-2pt]
         \scriptsize rubric mean};
  \draw[flow] (judge.west) -- (board.east);

  % gold answers feed the judge, and visibly trace back to the SAME graph
  \draw[flow] (gold.east) to[out=0,in=205]
        node[pos=0.5, below=0.3mm, font=\scriptsize, text=black!60] {gold}
        (judge.west);
  % the dashed "scored against" link: gold (graph-projected) <-> judged answer
  \draw[scored] (graph.south) to[out=-90,in=180]
        node[pos=0.7, below=0.5mm, accent, font=\scriptsize\itshape]
        {scored against graph truth} (board.west);

\end{tikzpicture}

% Keep the existing \caption{...}\label{fig:teaser}; do not caption inside here.
% =============================================================================
\begin{tikzpicture}[
    font=\small,
    node distance=6mm,
    accent/.style={orange!75!red},
    % --- the bi-temporal source-of-truth graph (origin) ---
    graph/.style={cylinder, shape border rotate=90, aspect=0.32, draw=orange!75!red,
                  very thick, fill=orange!8, minimum width=20mm, minimum height=15mm,
                  align=center, inner sep=2pt, text=black},
    % --- generic stage box ---
    box/.style={draw=black!55, thick, rounded corners=2pt, fill=black!2,
                align=center, inner sep=4pt, minimum height=9mm, text width=24mm},
    % --- system-under-test chip ---
    sys/.style={draw=black!45, semithick, rounded corners=1.5pt, fill=white,
                font=\scriptsize\ttfamily, inner sep=2.5pt, minimum width=23mm,
                align=center},
    % --- competency tag ---
    tag/.style={draw=orange!70!red, semithick, rounded corners=1pt, fill=orange!10,
                font=\scriptsize, inner sep=1.6pt, text=black!85},
    % --- judge / output emphasis ---
    judge/.style={draw=orange!75!red, very thick, rounded corners=2pt, fill=orange!6,
                  align=center, inner sep=4pt, minimum height=9mm, text width=26mm},
    board/.style={draw=black!60, thick, rounded corners=2pt, fill=black!4,
                  align=center, inner sep=4pt, minimum height=9mm, text width=22mm},
    flow/.style={-{Latex[length=2mm,width=1.6mm]}, black!65, semithick},
    proj/.style={-{Latex[length=2mm,width=1.6mm]}, orange!75!red, thick},
    scored/.style={{Latex[length=1.6mm]}-{Latex[length=1.6mm]}, densely dashed,
                   orange!75!red, semithick},
  ]

  % ---- ORIGIN: bi-temporal source-of-truth graph -------------------------
  \node[graph] (graph) {\textbf{Source-of-truth}\\\textbf{graph}\\[1pt]
        \scriptsize\itshape bi-temporal\\[-1pt]
        \scriptsize\ttfamily valid/invalid\\[-2pt]
        \scriptsize\ttfamily ingested/expired};
  \node[font=\scriptsize\itshape, text=black!55, above=1mm of graph]
        {synthetic Helix org};

  % ---- DUAL PROJECTION targets (the key idea) ----------------------------
  \node[box, above right=6mm and 16mm of graph, text width=27mm] (corpus)
        {\textbf{Multi-channel corpus}\\[1pt]
         \scriptsize email $\cdot$ Slack $\cdot$ transcripts\\[-2pt]
         \scriptsize tickets $\cdot$ docs $\cdot$ CRM};
  \node[box, below right=6mm and 16mm of graph, text width=27mm] (gold)
        {\textbf{Questions + gold}\\[1pt]
         \scriptsize competency-tagged Qs\\[-2pt]
         \scriptsize deterministic gold answers};

  % competency tags C1-C6 hanging under the questions
  \node[tag, below=2.5mm of gold] (tags)
        {C1 supers.\ $\cdot$ C2 prov.\ $\cdot$ C3 bi-temp.\ $\cdot$
         C4 audit $\cdot$ C5 justif.\ $\cdot$ C6 contra.};

  % projection arrows from the one graph (orange = "projected from graph")
  \draw[proj] (graph.east) to[out=35,in=180]
        node[pos=0.55, above=-0.5mm, accent, font=\scriptsize\itshape] {project}
        (corpus.west);
  \draw[proj] (graph.east) to[out=-35,in=180] (gold.west);

  % ---- SYSTEMS UNDER TEST -------------------------------------------------
  \node[box, right=14mm of corpus, text width=15mm, fill=white, draw=black!45]
        (ingest) {\textbf{ingest}\\\scriptsize+ retrieve};
  \node[sys, above right=2mm and 9mm of ingest] (s1) {gbrain};
  \node[sys, below=1.2mm of s1] (s2) {zep-cloud};
  \node[sys, below=1.2mm of s2] (s3) {mem0-platform};
  \node[sys, below=1.2mm of s3] (s4) {graphify-oss};
  \begin{scope}[on background layer]
    \node[draw=black!30, rounded corners=2pt, fill=black!2, fit=(s1)(s4),
          inner sep=2.5pt] (systems) {};
  \end{scope}
  \node[font=\scriptsize\itshape, text=black!55, above=0.5mm of systems]
        {memory systems under test};

  \draw[flow] (corpus.east) -- (ingest.west);
  \draw[flow] (ingest.east) -- (systems.west);

  % ---- ANSWERER -----------------------------------------------------------
  \node[box, right=10mm of systems, text width=20mm] (answerer)
        {\textbf{Answerer}\\[1pt]
         \scriptsize shared neutral\\[-2pt]
         \scriptsize (or native synth.)};
  \draw[flow] (systems.east) -- (answerer.west);

  % ---- JUDGE (capability-aware, rubric) ----------------------------------
  \node[judge, below=12mm of answerer] (judge)
        {\textbf{Capability-aware judge}\\[1pt]
         \scriptsize rubric facet coverage\\[-2pt]
         \scriptsize readiness-gated};
  \draw[flow] (answerer.south) -- node[right=0.5mm, font=\scriptsize, text=black!60]
        {answer} (judge.north);

  % ---- LEADERBOARD --------------------------------------------------------
  \node[board, left=12mm of judge] (board)
        {\textbf{Leaderboard}\\[1pt]
         \scriptsize per-category\\[-2pt]
         \scriptsize rubric mean};
  \draw[flow] (judge.west) -- (board.east);

  % gold answers feed the judge, and visibly trace back to the SAME graph
  \draw[flow] (gold.east) to[out=0,in=205]
        node[pos=0.5, below=0.3mm, font=\scriptsize, text=black!60] {gold}
        (judge.west);
  % the dashed "scored against" link: gold (graph-projected) <-> judged answer
  \draw[scored] (graph.south) to[out=-90,in=180]
        node[pos=0.7, below=0.5mm, accent, font=\scriptsize\itshape]
        {scored against graph truth} (board.west);

\end{tikzpicture}
}
  \caption{\textbf{The \bench{} evaluation pipeline.} A multi-channel, multi-author
    synthetic organizational corpus (email, Slack, meeting transcripts, tickets,
    wikis, CRM) is ingested by each memory system. Competency-tagged questions
    (the six small/medium-tier competencies: supersession, decision provenance,
    bi-temporal, audit replay, justification chains, contradiction) are answered
    through a shared answerer
    and scored by a fixed LLM-as-judge rubric, with a readiness gate ensuring hosted
    asynchronous systems are fully ingested before querying. The strongest public
    system answers only $\sim$$39\%$ of the task.}
  \label{fig:teaser}
\end{figure}
